\section{有符号整数除法：符号归一化、迭代与一次提交}

恢复余数除法已经给出了无符号核心：移位寄存器保存部分余数和待生成的商，每轮试减一次除数，再按结果决定当前商位。处理器要执行有符号除法，还需要在这条核心路径的前后各加一层控制。入口层识别符号和异常，把操作数转换成可供无符号核心处理的幅值；出口层恢复商与余数的符号，并在同一个体系结构提交点写回结果或报告异常。

以采用\hwkey{向零截断}（truncate toward zero）的常见整数 ISA 语义为例，商的小数部分直接舍去，余数满足
\[
    \text{dividend}=\text{quotient}\times\text{divisor}+\text{remainder},
    \qquad |\text{remainder}|<|\text{divisor}|.
\]
非零余数与被除数同号。这个规则不是排版约定，而是出口符号校正逻辑必须实现的接口契约。

\topic{案例：用同一个无符号核心计算 \texorpdfstring{$-13\div3$}{-13/3}}

设数据宽度为 8 bit，被除数为 \code{11110011}（$-13$），除数为 \code{00000011}（$+3$）。入口先锁存原始符号：被除数符号为 1，除数符号为 0，因此商符号为二者异或，即 1；余数符号直接继承被除数符号。随后把两个操作数转换为幅值 13 和 3，交给无符号迭代器。核心得到幅值商 4、幅值余数 1，出口分别施加预先保存的符号，形成商 $-4$、余数 $-1$。

\begin{longtable}{L{0.17\linewidth}L{0.24\linewidth}L{0.27\linewidth}L{0.22\linewidth}}
\toprule
阶段 & 保存的状态 & 数据路径动作 & 此时能否写回 \\
\midrule
入口检查 & 原始操作数、两个符号位 & 检查除零和唯一溢出组合 & 不能，异常尚未定案 \\
幅值归一化 & 商符号、余数符号 & 负数取补，形成 13 与 3 & 不能，只有内部幅值 \\
无符号迭代 & 部分余数、部分商、轮次计数 & 逐轮移位、试减并生成商位 & 不能，中间商尚不完整 \\
符号恢复 & 幅值商 4、幅值余数 1 & 商取负、余数取负 & 不能，仍要检查最终状态 \\
体系结构提交 & 商 $-4$、余数 $-1$、异常为空 & 一次写入目标寄存器与状态 & 可以，满足 $-13=(-4)\times3+(-1)$ \\
\bottomrule
\end{longtable}

这里有两个容易混淆的符号。商符号由“被除数符号异或除数符号”决定；余数符号不做异或，而是跟随被除数。若误让余数跟随商，$-13\div3$ 会得到商 $-4$、余数 $+1$，代回恒等式变成 $-11$，接口契约立即被破坏。

\begin{pdfdiagram}[x=1cm,y=1cm]
  \node[hw state, minimum width=2.4cm] (capture) at (1.3,3.8) {锁存操作数\\符号与原始位型};
  \node[hw control, minimum width=2.5cm] (check) at (4.3,3.8) {入口检查\\除零/溢出};
  \node[hw compute, minimum width=2.5cm] (magnitude) at (7.4,3.8) {幅值归一化\\条件取补};
  \node[hw compute, minimum width=2.7cm] (iterate) at (10.7,3.8) {无符号迭代\\余数/商/轮次};
  \node[hw compute, minimum width=2.7cm] (signfix) at (10.7,1.2) {符号恢复\\商与余数分别处理};
  \node[hw state, minimum width=2.7cm] (commit) at (7.4,1.2) {一次提交\\写回或精确异常};
  \node[hw control, minimum width=2.5cm] (fault) at (4.3,1.2) {异常状态\\不启动迭代器};
  \draw[hw bus] (capture) -- (check);
  \draw[hw bus] (check) -- node[above, hw label]{合法} (magnitude);
  \draw[hw bus] (magnitude) -- (iterate);
  \draw[hw bus] (iterate) -- (signfix);
  \draw[hw bus] (signfix) -- (commit);
  \draw[hw return] (check) -- node[right, hw label]{除零或溢出} (fault);
  \draw[hw return] (fault) -- (commit);
\end{pdfdiagram}
\diagramnote{有符号整数除法的控制闭环。上排只在入口合法时启动耗时的无符号迭代；下排把正常结果或异常送到同一个提交点。中间余数和部分商均是微架构状态，不能提前改变体系结构目标寄存器。}

\topic{除零与最小负数除以负一必须在入口闭合}

除数为 0 时，没有任何商余数对能够满足除法恒等式。硬件在启动迭代前比较除数是否为零，可以避免迭代器做完全部轮次才发现失败。体系结构可以规定产生陷阱、返回固定值或设置状态位；无论选择哪种语义，目标寄存器和异常状态必须在同一个精确边界更新。

二进制补码还有一个唯一的有符号溢出组合。$n$ bit 的最小负数是 $-2^{n-1}$，而最大正数只有 $2^{n-1}-1$。因此
\[
    -2^{n-1}\div(-1)=+2^{n-1}
\]
无法放回同宽目标寄存器。8 bit 示例是 $-128\div(-1)=+128$。入口可以直接比较“被除数等于 \code{10000000} 且除数等于 \code{11111111}”，也可以让内部幅值路径临时扩成 $n+1$ bit，再在出口检测结果越界。前者省去无意义迭代，后者便于复用更宽的数据通路；两种实现都必须给出同一体系结构结果。

\begin{longtable}{L{0.20\linewidth}L{0.25\linewidth}L{0.28\linewidth}L{0.17\linewidth}}
\toprule
输入类别 & 入口决策 & 可见结果 & 为什么不进入普通迭代 \\
\midrule
除数为 0 & 形成除零条件 & 按 ISA 产生陷阱、状态或规定结果 & 数学上不存在合法商余数对 \\
被除数为 0、除数非零 & 直接形成 0 与 0 & 商 0、余数 0 & 无需生成任何非零商位 \\
除数为 $+1$ 或 $-1$ & 直通或条件取补 & 商由符号规则形成，余数为 0 & 幅值结果已知 \\
\code{INT\_MIN}/$-1$ & 形成有符号溢出 & 按 ISA 报告溢出或规定结果 & 正结果超出同宽补码范围 \\
一般输入 & 启动迭代器 & 完整商余数或后续检测结果 & 每一位商依赖上一轮余数 \\
\bottomrule
\end{longtable}

\topic{可变时延只能改变完成时间，不能改变提交语义}

固定轮次恢复余数法对 $n$ bit 商执行 $n$ 轮，完成时间容易预测。真实除法器还可以利用前导零计数跳过必为零的高位商，识别小除数或特殊操作数，也可以采用 SRT、基于倒数近似的乘法迭代或分段流水结构。这些优化使不同输入的延迟不同，但发射端仍要为该操作分配状态，相关消费者仍要等待完成标志，异常仍要与原指令身份一起到达提交点。

若中断或页错误发生在除法迭代期间，简单多周期处理器可以保存迭代器状态，或丢弃内部状态并在返回后从指令起点重做；乱序处理器通常让除法器继续或取消该操作，但只有对应指令到达退休位置时才提交结果。恢复策略可以不同，体系结构观察到的效果必须相同：要么整条除法完成一次，要么目标寄存器保持旧值并在精确边界进入异常处理。
